\section{Libraries and public extension interfaces}
\subsection{Full loading or the minimal mode}
\key{\usepackage{gradbars}} retains full loading and previous public commands.
Alternatively, list all feature libraries explicitly:
\begin{source}
\usepackage[libraries={stack,spark,heat,csv,table}]{gradbars}
\end{source}
This is equivalent to default full loading; shorten the list to select fewer libraries.
Select only the shared engine, then load features in the preamble:
\begin{source}
\usepackage[libraries=none]{gradbars}
\usegradbarslibrary{stack,spark}
\end{source}
The package option \key{libraries={heat,spark}} is also available. Core contains
ordinary bars, comparisons/ranges, references, quality controls and legends.
\key{stack}, \key{spark}, \key{heat} and \key{csv} add their respective engines;
\key{table} adds native visual tables and all four dependencies. Repeated loading
is safe. Unknown libraries, missing files, cycles and unloaded frontends produce
clear errors. Load libraries in the preamble; siunitx v3 remains explicitly optional.
\subsection{Complete list of built-in libraries}
Version 1.0.0 registers six built-in library names. Core is always loaded;
the other five provide optional features:
\begin{description}
\item[\key{core}] Basic bars, comparisons, dumbbells, ranges, reference marks,
quality controls and legends, plus number handling, settings and accessibility
tagging. Always loaded automatically; no explicit selection is needed.
\item[\key{stack}] Signed stacked bars, segment labels and stacked legends.
The main command is \key{\gradstack}; depends on core.
\item[\key{spark}] Miniature line, bar and win-loss series.
The main command is \key{\gradspark}; depends on core.
\item[\key{heat}] Sequential and diverging heat cells with a configurable center.
The main command is \key{\gradheat}; depends on core.
\item[\key{csv}] CSV import, cell access, CSV tables and collected columns in
existing tables. Main commands include \key{\gradbarsloadcsv} and
\key{\gradbarscolumn}; depends on core.
\item[\key{table}] Native \key{gradtable} tables with columns, groups, summaries,
shared scales, ranking, styles and pagination. Automatically loads core,
\key{stack}, \key{spark}, \key{heat} and \key{csv}.
\end{description}
List all six names as follows. Explicitly selecting core is valid but unnecessary:
\begin{source}
\usepackage[libraries={core,stack,spark,heat,csv,table}]{gradbars}
\end{source}
Alternatively, start in minimal mode and load the libraries in the preamble:
\begin{source}
\usepackage[libraries=none]{gradbars}
\usegradbarslibrary{core,stack,spark,heat,csv,table}
\end{source}
The values \key{all} and \key{none} are modes of the package option
\key{libraries}, not library names. Select \key{libraries=table} for native
tables; the loader supplies all dependencies. The files named \key{keys},
\key{numbers}, \key{render-tikz}, \key{tagging}, \key{api}, \key{table-data}
and \key{table-layout}, as well as \key{model}, are internal modules, not separately loadable public libraries.
Register contributed libraries with \key{\gradbarsdeclarelibrary} first.
The \key{demo} library in this chapter is a companion example, not built in.

\subsection{Unified settings}
\key{\gradbarsset} and \key{\gradbarssetup} share options, grouping and precedence.
Local settings and per-graphic overrides retain their previous scope.
\begin{demo}{Settings with a local override}
\gradbarsset{max=100,precision=2}
\gradbar{72.125}
{\gradbarsset{color=teal}\gradbar{50}}
\end{demo}
\subsection{Stable registration callbacks}
These public interfaces follow the version 1 stability contract:
\begin{itemize}
\item Graphics: \key{\gradbarsdeclaretype} and \key{\gradgraphic}.
\item Formatting: \key{\gradbarsdeclareformatter} and \key{\gradbarsformatnumber}.
\item Libraries: \key{\gradbarsdeclarelibrary}.
\end{itemize}
A scalar callback receives one validated value or \key{NA}
and should delegate to public graphics. A formatter receives display value,
precision and number-format mode; return just the number, with the engine adding
unit/percent suffixes. Registration is global and duplicate names are errors.
Names start with a lower-case letter and use letters, digits and hyphens.
\begin{source}
\gradbarsdeclaretype{score}{\gradbar[color=teal]{#1}}
\gradbarsdeclareformatter{boxed}{%
  \fbox{\gradbarsformatnumber[precision=#2,number format=#3]{#1}}}
\gradgraphic{score}{72.125}
\gradbar[formatter=boxed]{72.125}
\end{source}
Custom scalar types do not add new native \key{gradcolumn} types. Custom formatters
work in graphics and ordinary number columns, but do not support
\key{number align=decimal}. Data calculations never consume formatted strings.
\subsection{A public-only contributed library}
The companion \key{docs/libraries/gradbars-contrib-demo.sty} registers a renderer,
formatter and table style. Register its dependencies and load it in the preamble:
\begin{source}
\gradbarsdeclarelibrary{demo}{table}{docs/libraries/gradbars-contrib-demo.sty}
\usegradbarslibrary{demo}
\end{source}
\begin{demo}{Renderer, formatter and reusable table style}
\gradgraphic{demo-score}{72.125}\par
\gradbar[formatter=demo-boxed]{72.125}\par
\begin{center}
\begin{gradtable}[align=c,header align=c,style=demo-report,width=\linewidth]
  \gradcolumn{name}{Method}
  \gradcolumn[type=number,formatter=demo-boxed,aggregate=mean]{value}{Score}
  \gradrow{A,72}\gradrow{B,90}\gradsummary{Mean}
\end{gradtable}
\end{center}
\end{demo}
The example uses no private variables. See \key{docs/developer.md} and
\key{docs/stability.md} for callback limits, dependencies and compatibility.
\subsection{Install the complete release}
For a CTAN source archive, first run \key{texlua scripts/unpack.lua} from its root
to create \key{sty/} and extract the runtime files. Install all matching
\key{.sty} files from \key{sty/} together, including \key{tblrlibgradbars.sty}; replacing only the entry file is insufficient.
Native versions must match. In the GitHub checkout, use \key{l3build check},
\key{l3build doc}, \key{l3build ctan} and the detailed
\key{texlua scripts/check.lua} regressions. The CTAN source archive omits
\key{tests/}; it supports extraction and documentation rebuilds, while full
regression checks require the development checkout.
\subsection{Version 1 compatibility contract}
The documented public commands, keys, defaults and dependency rules are stable.
Compatible updates preserve existing interfaces, defaults and data semantics.
Removals or changes to defaults/data
semantics require a major release. Deprecated interfaces remain available
throughout version 1. Internal module files and private variables are excluded.
See \key{docs/stability.md} for scope, precedence and callback restrictions.
The supported baseline is LaTeX 2025-11-01, expl3 2025-12-29 and PGF/\manualTikZ{}
3.1.11a or newer. Install all twenty-three matching runtime files, including optional feature libraries and adapters, together.
The generated \key{gradbars.dtx} and \key{gradbars.ins} reproduce those modules
using docstrip; \key{texlua scripts/source.lua --check} verifies consistency.
The l3doc source guide complements the illustrated manuals.

\subsection{Declaring convenient graphic commands}
\key{\gradbarsnewcommand[declaration and graphic keys]{\command}} creates a scoped
command without replacing existing names. Defaults are bar, one mandatory argument,
and data from the first argument. Only \key{m} and \key{O{default}} argument specifications are
accepted, with one to nine arguments.
\begin{demo}{A reusable command with trailing optional overrides}
\begingroup
\gradbarsnewcommand[type=bar,args={m O{}},data={#1},options={#2},
 empty=skip,theme=gray,min=0,max=1,width=12mm,height=5pt,
 precision=3,label position=right,label distance=1mm]{\servicebar}
\servicebar{0.960}\par\medskip
\servicebar{0.720}[color=teal,width=20mm]\par\medskip
\servicebar{0.400}[label position=above,label reference=value,
 label distance=2pt,label xshift=-.3em,label yshift=1pt]
\endgroup
\end{demo}
Use \key{args={m m}} for two mandatory arguments, mapping the first to data
and the second to options. \key{args={O{} m}} places optional settings first;
map the second argument to data and the first to options.
Supported types are bar, heat, stack, spark, compare, dumbbell, range, boxplot,
and registered scalar or pair types. Pair types require two grouped values;
see Section~\ref{sec:pair-extensions-en} for their independent record contract:
\begin{source}
\gradbarsnewcommand[type=dumbbell,args={m m},
  data={{#1}{#2}}]{\pairbar}
\end{source}
Required libraries remain explicit.
Whitespace-only input follows \key{empty=missing|skip|error}; zero is always data,
and explicit NA stays missing. Options apply after surrounding setup and fixed
command keys, without changing named-scale or allocated-table-width rules.

\subsection{Color and explicit label placement}
Choose a bar theme, a solid \key{color=gray!60}, or explicit
\key{fill=gradient}, \key{left color=blue} and \key{right color=cyan!20}. Later settings override
earlier ones. Heat continues to use heat palette/low/high color keys.
\key{label position} accepts above, below, left, right, inside or none, opting into native bar/Heat
main-label placement. Other types retain their own label contracts and diagnose
unsupported explicit placement. \key{label reference=track} uses the full track;
bars also accept value, using the clipped endpoint or track center when missing.
Heat requires track. \key{label distance} defaults to .6em, must be nonnegative,
and is ignored inside. \key{label xshift} and \key{label yshift} default to 0pt,
with positive values moving right and up. Lengths are signed decimal numbers with
pt, mm, cm, in, bp, pc, dd, cc, sp, em or ex; em/ex use the actual label font.
Unitless values and general expressions are rejected.
\key{label anchor=auto} selects the appropriate edge; explicit center, north,
south, east, west and the four two-word corner anchors are also supported.
Actual shifted bounds contribute width, height and depth; fixed cells still need
enough space. No automatic collision avoidance is promised. Later legacy
\key{label=outside} etc. restores old placement. Explicit placement requires
native TikZ or \key{backend fallback=tikz}; schema 1 is unchanged. Area contrast
is outside this contract; automatic text color is not a contrast certification.

Tagged tabularray examples use explicit column widths and opt-in table tagging; natural-width tagged layout has not been verified.
